% !TEX program = xelatex
% 编译方式：在本目录执行两遍 xelatex（详见 docs/_manual_common/README.md）。
% 源码依据：src/api/ 与 include/hacdcpf/api/（hacdcpf.hpp / solver_capabilities.hpp）及 docs/reference/runtime_api.md。
\documentclass[UTF8,fontset=fandol,a4paper,11pt]{ctexart}
\usepackage{../../_manual_common/hysim_manual}

\title{\textbf{交直流混合高弹性能源电力系统仿真平台}\\[0.5em]
       {\Large 公共 API 门面技术手册}\\[0.3em]
       {\large 求解入口族、能力门控、异常语义与工业级评估}}
\author{HySim-XJTU-HRPES（西安交通大学 HRPES 团队）}
\date{\today}

\begin{document}
\maketitle
\begin{abstract}
\noindent 本手册依据 HySim-XJTU-HRPES 平台公共 API 门面
（\srcpath{include/hacdcpf/api/hacdcpf.hpp}、\srcpath{solver\_capabilities.hpp}，
实现于 \srcpath{src/api/hacdcpf.cpp}）整理。门面在命名空间 \texttt{hacdcpf} 下以稳定
自由函数聚合全库的潮流（PF）、最优潮流（OPF）、动态、碳流、市场、可靠性、弹性、
时序与校验子系统，并经 \srcpath{get\_solver\_capabilities()} 在运行期报告后端可用性。
可选后端由能力标志门控，缺失后端如实报告不可用而非静默降级；\srcpath{verify\_opf\_result}
为独立的解后审计，不改变求解本身。GUI/自动化 HTTP 面属独立层
（\srcpath{tests/run\_gui\_server.cpp}），不在 \srcpath{src/api}。全部结论以源码为准；版式
与《潮流计算技术手册》《最优潮流技术手册》一致。命名空间为 \texttt{hacdcpf}。
\end{abstract}

\tableofcontents
\newpage

\section{阅读指南与约定}
\subsection{数据来源}
\begin{itemize}
  \item \srcpath{include/hacdcpf/api/hacdcpf.hpp} —— 门面自由函数与会话/句柄类型；
  \item \srcpath{include/hacdcpf/api/solver\_capabilities.hpp} —— \texttt{SolverCapabilities}
        与运行期能力查询；
  \item \srcpath{src/api/hacdcpf.cpp} —— 门面与能力查询的实现；
  \item \srcpath{docs/reference/runtime\_api.md} —— HTTP/GUI \texttt{/api/v1} 服务契约（位于
        \srcpath{tests/run\_gui\_server.cpp}）；
  \item \srcpath{docs/reference/python\_api.md} —— Python 绑定面；
  \item \srcpath{tests/test\_hacdcpf.cpp}、\srcpath{tests/test\_solver\_capabilities.cpp}
        —— 验证依据。
\end{itemize}

\subsection{单位制与索引空间}
门面不改变各子系统的单位与索引口径：潮流量以 MW/Mvar/标幺电压表达，结果向量的索引
空间与对应求解器一致（详见各专项手册）。组件对外以稳定 \texttt{.index} 标识；AC/DC 域
分开，门面不引入新的跨域裸 ID。所有返回结构沿用子系统既有的有效性/诚实字段。

\subsection{符号与命名约定}
门面函数名以子系统语义前缀区分（\srcpath{solve\_power\_flow\_*}、\srcpath{solve\_*\_opf}、
\srcpath{run\_*}、\srcpath{verify\_*}、\srcpath{validate\_*}）；\texttt{safe\_} 前缀表示
返回 \texttt{Result<T>} 且可选先校验；能力标志以 \fld{has\_*}（后端存在）与
\fld{supports\_*}（特性可用）两族命名。

\subsection{诚实口径}
门面是薄封装：能力标志如实反映\textbf{编译期}后端可用性，运行期查询不会“凭空”启用缺失
后端；\srcpath{solve\_power\_flow\_helm} 仅支持 AC（混合 AC/DC 与电压相关 ZIP 负荷被显式
拒绝并给出诊断）；\srcpath{safe\_solve\_power\_flow} 在 \fld{validate\_input}=true 时先校验再
求解，而抛异常版 \srcpath{solve\_power\_flow} 不校验；\srcpath{verify\_opf\_result} 是独立的
解后审计，不改变求解结果本身。

\section{模块定位与工程应用场景}
\subsection{功能定位}
API 门面是全平台\textbf{唯一稳定的对外调用面}：它对投影、装配、符号分析、求解器工厂等
内部细节做封装，把 PF/OPF/动态/碳/市场/可靠性/弹性/时序/校验各子系统的入口以命名空间
\texttt{hacdcpf} 下的自由函数再导出，并以 \texttt{SolverCapabilities} 报告可选后端。它是
GUI 后端、Python 绑定、批处理工具与回归测试共同依赖的接口边界。

\subsection{工程应用场景}
典型场景：GUI/HTTP 后端按 \texttt{method} 分派到门面求解入口；Python 客户端经绑定调用
\srcpath{solve\_power\_flow}/\srcpath{solve\_ac\_opf} 做批量算例；时序生产模拟经
\texttt{PreparedPowerFlowSession}/\texttt{SolverHandle} 在时间轴上复用符号分解；工具链在
运行期查询 \srcpath{get\_solver\_capabilities()} 以决定是否启用 MILP/ETAP/OpenDSS 路径。

\section{输入、输出与边界条件}
\subsection{输入}
已装配的 \texttt{HybridPowerSystem} 与对应子系统的选项结构体（\texttt{PowerFlowOptions}、
\texttt{OPFOptions} 等）；\texttt{safe\_} 变体额外接受 \fld{validate\_input} 开关；句柄/会话
接口接受兼容性快照签名以判定复用。

\subsection{输出}
各子系统既有结果结构（\texttt{PowerFlowResult}、\texttt{OPFResult} 等）；\texttt{safe\_} 变体
返回 \texttt{Result<T>}（成功值或错误）；\srcpath{get\_solver\_capabilities()} 返回
\texttt{SolverCapabilities}；\srcpath{verify\_opf\_result} 把审计写入 \fld{result.audit}。

\subsection{边界条件}
门面不校验（除 \texttt{safe\_} 变体）；HELM 仅 AC；句柄复用要求拓扑与后端签名兼容
（否则重建）；能力门控特性在后端缺失时相应入口返回不可用错误而非降级近似。

\section{变量、参数与符号定义}
\subsection{关键数据结构}
\compmeta{SolverCapabilities}{include/hacdcpf/api/solver\_capabilities.hpp}
线性代数后端 \fld{has\_sparse\_lu}（恒 true）/\fld{has\_klu}/\fld{has\_umfpack}/\fld{has\_accelerate}；
MIP/NLP 后端 \fld{has\_highs}/\fld{has\_gurobi}/\fld{has\_ipopt}/\fld{has\_scip}/\fld{has\_native\_ipm}
（恒 true）；I/O 能力 \fld{has\_excel}/\fld{has\_opendss}；特性标志
\fld{supports\_integer\_variables}、\fld{supports\_ac\_opf}（true）、\fld{supports\_dc\_opf}（true）、
\fld{supports\_three\_phase}（true）。

\compmeta{PreparedPowerFlowSession}{include/hacdcpf/api/hacdcpf.hpp}
持有型重复求解上下文，成员 \srcpath{solve(sys)}、\srcpath{reset()}；move-only（拷贝已删除），
采用 pimpl \texttt{Impl}。另有前置声明的不透明句柄 \texttt{SolverHandle}，提供 C 风格可复用
PF/DC 句柄。

\subsection{参数字典}
下表列出 \texttt{SolverCapabilities} 的关键能力标志（“默认值”列“编译期”表示由构建
配置决定，“派生”表示由其他标志推导）。
\begin{paramtable}
\fld{has\_sparse\_lu} & — & — & true/false & true & 内置稀疏 LU（恒可用）\\
\fld{has\_klu} & — & — & true/false & 编译期 & KLU 稀疏后端\\
\fld{has\_umfpack} & — & — & true/false & 编译期 & UMFPACK 稀疏后端\\
\fld{has\_accelerate} & — & — & true/false & 编译期 & Apple Accelerate 稀疏后端\\
\fld{has\_highs} & — & — & true/false & 编译期 & HiGHS LP/MILP 后端\\
\fld{has\_gurobi} & — & — & true/false & 编译期 & Gurobi 后端\\
\fld{has\_ipopt} & — & — & true/false & 编译期 & Ipopt NLP 后端\\
\fld{has\_scip} & — & — & true/false & 编译期 & SCIP 后端\\
\fld{has\_native\_ipm} & — & — & true/false & true & 内置原生 IPM（恒可用）\\
\fld{has\_excel} & — & — & true/false & 编译期 & OpenXLSX/ETAP I/O 能力\\
\fld{has\_opendss} & — & — & true/false & 编译期 & OpenDSS 桥接能力\\
\fld{supports\_integer\_variables} & — & — & true/false & 派生 & 需 HiGHS/Gurobi/SCIP 之一\\
\fld{supports\_ac\_opf} & — & — & true/false & true & AC OPF 支持\\
\fld{supports\_dc\_opf} & — & — & true/false & true & DC OPF 支持\\
\fld{supports\_three\_phase} & — & — & true/false & true & 三相支持\\
\end{paramtable}

\section{数学模型与约束条件：能力门控}
门面本身不含数学模型，其“模型”是\textbf{能力门控的布尔代数}。可选特性按下述逻辑
门控（$\vee$ 为逻辑或）：
\[
\texttt{supports\_integer\_variables}\;\Longleftrightarrow\;
\texttt{has\_highs}\ \vee\ \texttt{has\_gurobi}\ \vee\ \texttt{has\_scip}.
\]
ETAP I/O 可用当且仅当 \fld{has\_excel} 为真；OpenDSS I/O 可用当且仅当 \fld{has\_opendss}
为真。约束：\fld{has\_sparse\_lu} 与 \fld{has\_native\_ipm} 恒为 true（内置），故稀疏直接解
与原生 IPM 在任意构型下均可用；其余后端缺失时相应入口如实报告不可用而非静默降级。

\section{算法流程：门面分派}
\begin{figure}[H]\centering
\begin{tikzpicture}[font=\footnotesize,
  box/.style={draw,rounded corners,minimum height=0.85cm,align=center,inner sep=3pt},arr/.style={->,thick}]
  \node[box](call) at (0,0){调用方\\GUI/Python/工具};
  \node[box](fac) at (3.4,0){门面自由函数\\hacdcpf::solve\_*};
  \node[box](cap) at (7.0,0){能力门控\\get\_solver\_capabilities};
  \node[box](sub) at (10.4,0){子系统求解器\\PF/OPF/…};
  \node[box](sess) at (3.4,-1.9){会话/句柄\\符号分解复用};
  \node[box](res) at (7.0,-1.9){结果结构\\子系统既有口径};
  \node[box](aud) at (10.4,-1.9){解后审计(可选)\\verify\_opf\_result};
  \draw[arr](call)--(fac);\draw[arr](fac)--(cap);\draw[arr](cap)--(sub);
  \draw[arr](fac.south)|-(sess.east);\draw[arr](sub)--(res);\draw[arr](res)--(aud);
\end{tikzpicture}
\caption{门面分派与能力门控（会话/句柄在兼容快照间复用符号分析）}
\end{figure}
调用方无需了解投影、装配与符号分析细节；\texttt{PreparedPowerFlowSession} 与
\texttt{SolverHandle} 在兼容快照之间复用投影/装配/符号分析结果，摊薄重复求解开销。

\section{工程假设与数值稳定性分析}
\subsection{工程假设}
门面假设子系统各自维护其数值与收敛口径，门面层只做转发与门控，不叠加额外近似；假设
能力标志在进程生命周期内不变（编译期决定）。

\subsection{数值稳定性}
门面不引入新的数值误差；数值稳定性完全由被调子系统决定（见各专项手册）。会话/句柄复用
符号分解时，若拓扑签名不匹配则强制重建，避免使用过期的符号模式导致错误分解。

\section{复杂度分析}
门面转发为 $O(1)$ 额外开销；能力查询为常数时间的标志读取。真正的计算复杂度由被调求解器
决定。会话/句柄的价值在于把符号分析成本（一次 $O(\text{nnz})$ 的 \texttt{analyzePattern}）
摊薄到时间轴上的多次数值分解，时序生产模拟因此显著提速。

\section{异常工况处理}
\begin{itemize}
  \item \textbf{输入非法}：\texttt{safe\_} 变体在 \fld{validate\_input}=true 时返回错误
        \texttt{Result}；抛异常版直接抛出，调用方需自行捕获；
  \item \textbf{后端缺失}：门控特性入口返回不可用错误，不静默降级为近似方法；
  \item \textbf{HELM 遇混合/ZIP}：显式拒绝并诊断，不尝试不支持的求解；
  \item \textbf{句柄/会话签名不兼容}：自动重建符号分解，保证正确性优先于复用；
  \item \textbf{求解不收敛}：透传子系统的收敛标志与诊断，门面不掩盖失败。
\end{itemize}

\section{与其他模块的接口关系}
门面头文件汇入 PF/OPF/动态/碳流/市场/可靠性/弹性/时序/校验/EV/综合能源各子系统，是它们
对外的统一出口；GUI 后端（\srcpath{tests/run\_gui\_server.cpp}）与 Python 绑定
（\srcpath{docs/reference/python\_api.md}）均建立在门面之上；\srcpath{verify\_opf\_result} 与 OPF
子系统解耦，作为独立审计层。门面不含网络层，HTTP 面属独立可执行。

\section{测试与验证建议}
注册测试：\srcpath{test\_hacdcpf}、\srcpath{test\_solver\_capabilities}、\srcpath{test\_hacdcdss}，
覆盖门面入口、运行期能力查询与端到端装配。建议补充：能力门控真值表（各后端组合下
\fld{supports\_integer\_variables} 等派生标志）、\texttt{safe\_} 与抛异常版的行为对照、
HELM 拒绝混合/ZIP 的诊断断言、以及会话/句柄在拓扑变更时的重建正确性。

\section{工业级评估指标}
建议纳入：门面转发开销（应为纳秒级，可忽略）；会话/句柄复用带来的时序仿真加速比；
能力查询与实际后端行为的一致率（门控真值表 100\% 匹配）；\texttt{safe\_} 变体对非法输入
的拦截率与误报率；解后审计 \srcpath{verify\_opf\_result} 对已知不可行解的检出率。

\section{适用范围与局限性}
\begin{itemize}
  \item \srcpath{solve\_power\_flow\_helm} 仅支持 AC；混合 AC/DC 与电压相关 ZIP 负荷被
        显式拒绝并给出诊断；
  \item \srcpath{safe\_solve\_power\_flow} 在 \fld{validate\_input}=true 时先校验，抛异常版
        \srcpath{solve\_power\_flow} 不校验；
  \item 能力标志门控可选特性：\fld{supports\_integer\_variables} 需 HiGHS/Gurobi/SCIP
        之一，\fld{has\_excel}/\fld{has\_opendss} 门控 ETAP/OpenDSS I/O；
  \item \srcpath{verify\_opf\_result} 为独立解后可行性审计（功率平衡、电压/支路/机组
        限值、目标一致性），结果写入 \fld{result.audit}，与求解过程解耦；
  \item GUI/自动化 HTTP 面在 \srcpath{tests/run\_gui\_server.cpp}，非 \srcpath{src/api}；
        门面本身不含网络层。
\end{itemize}

\section{后续改进方向}
\begin{enumerate}
  \item 将求解器工厂层（\srcpath{solver\_interface.cpp}）与门面对接，令生产路径经统一抽象；
  \item 为 \texttt{safe\_} 族补齐全子系统覆盖，逐步收敛抛异常版为内部实现；
  \item 扩展 \texttt{SolverCapabilities} 报告运行期资源（线程数、内存上限）以支持自适应调度；
  \item 将 \srcpath{verify\_opf\_result} 式解后审计推广到 PF/动态等其他子系统。
\end{enumerate}

% theory_api_contract.tex -- 公共 API 理论层
\section{门面 API 的能力、生命周期与结果传播理论}\label{sec:api-theory}

\begin{theorynote}
API 门面不增加电气方程；本章形式化的是能力门控、异常边界、句柄复用和独立解后审计。
具体求解模型仍由 PF/OPF/动态等专项手册负责。
\end{theorynote}

\subsection{能力门控}
令 $H$、$G$、$S$ 分别表示 HiGHS、Gurobi、SCIP 可用，则整数变量能力为
\begin{equation}
supports_{int}=H\lor G\lor S.
\end{equation}
\fld{has\_*} 是编译/链接能力，\fld{supports\_*} 是由能力组合推出的特性；二者不应被
调用方当作求解成功保证。

\subsection{安全入口与异常语义}
\texttt{safe\_solve\_power\_flow} 的调用契约是
\begin{equation}
validate\_input\land(\text{validation has Error})\Rightarrow Result<T>::error,
\end{equation}
而抛异常版本透传非法输入或求解失败。门面不把后端缺失、非收敛和 fallback 改写为成功。

\subsection{会话复用}
对网络结构签名 $\sigma$，缓存只在 $\sigma_{new}=\sigma_{cached}$ 时复用符号分析；
结构变化必须重建。数值注入和设定可以在同一结构上刷新。该策略只减少装配/符号分析成本，
不改变结果向量索引空间。

\subsection{与实现的对应}
\begin{center}\footnotesize
\begin{longtable}{@{}L{7.0cm}L{8.2cm}@{}}
\toprule 理论条目 & 实现状态\\\midrule
\endfirsthead\toprule 理论条目 & 实现状态\\\midrule\endhead
能力派生布尔式 & \implfull{src/api/hacdcpf.cpp:get_solver_capabilities}\\
安全 PF 入口 & \implfull{src/api/hacdcpf.cpp:safe_solve_power_flow}\\
符号签名缓存 & \implfull{src/api/hacdcpf.cpp:hash_system_signature}\\
OPF 解后审计 & \implfull{src/api/hacdcpf.cpp:verify_opf_result}\\
跨进程会话持久化 & \implnone\\
\bottomrule
\end{longtable}
\end{center}

% numerical_cross_validation.tex -- API 数值与契约证据
\section{能力真值表、入口回归与验证边界}\label{sec:api-numerical}

macOS Release 的 \srcpath{test\_solver\_capabilities} 通过 17 cases/61 assertions，覆盖能力
结构、可选后端派生标志和门面查询；\srcpath{test\_hacdcpf} 提供公共入口装配回归。
测试重点不是求解器精度，而是“能力报告与实际编译配置一致”。

\begin{center}\small
\begin{tabular}{@{}L{6.5cm}L{5.0cm}@{}}
\toprule 契约 & 当前证据\\\midrule
\fld{has\_sparse\_lu}=true、\fld{has\_native\_ipm}=true & Release 能力测试通过\\
\fld{supports\_integer\_variables}=$H\lor G\lor S$ & 17 cases/61 assertions\\
safe 入口返回结构化 Result & 公共 API 回归覆盖\\
句柄拓扑改变后重建 & 需专项生命周期测试，当前未独立量化\\
\bottomrule
\end{tabular}
\end{center}

当前没有跨语言 ABI、异常吞吐性能或句柄长期时序 benchmark；本章不把门面测试扩展成各
专项求解器的数值认证，也不把能力为 true 解释成特定算例必然收敛。

\section{跨模块工业级技术评价基线}

本章规定本手册所述模块进入工程算例、批量研究或生产服务前必须共同满足的评价基线。
具体模型公式、变量、约束和专用算法以正文相应章节为准；本章不扩大源码已经实现的范围。

\subsection{输入、输出、边界条件与数据结构审计}
输入审计应逐项检查有限性、单位、上下界、枚举值和引用完整性。系统对象必须先按所需层级完成
校验；AC 与 DC 母线采用域限定映射，组件稳定 \texttt{index}、向量位置和临时图索引不得混用。
输出审计必须记录单位、索引空间、模型范围、收敛或可行状态、后端、容差、迭代/节点计数和告警。
空系统、空时域、孤岛、重复 ID、零容量、退化约束与不可用可选后端均属于边界测试集。

\subsection{工程假设与数值稳定性分析}
模型使用的平衡、线性化、稳态、独立性、概率分布、时间离散或网络等值假设必须与应用场景匹配。
数值评价至少包含变量缩放、绝对/相对残差、可行性违反、条件数或枢轴质量、步长/阻尼、迭代停滞
和随机种子复现性。若采用正则化、平滑、回退、松弛或近似，应同时报告触发条件、参数值、对主指标
的敏感性以及是否改变模型口径；不得用“已返回结果”替代收敛或有效性证明。

\subsection{复杂度与资源分析}
令 $n_x$、$n_c$、$n_t$、$n_s$、$|V|$、$|E|$ 分别表示变量、约束、时段、场景、节点和边数。
报告应分别给出建模、求解、后处理的时间与内存。图遍历通常为 $O(|V|+|E|)$；逐时段/逐场景
流程至少线性增长为 $O(n_t n_s)$ 次子问题调用；稠密线性代数最坏可达 $O(n_x^3)$，稀疏方法则应
报告结构非零数、填充率和因子分解峰值内存；MILP/NLP 不给出虚假的多项式最坏界，而报告节点数、
间隙、时限和实测扩展曲线。复杂度结论必须基于固定硬件、固定后端和可复现算例族。

\subsection{异常工况处理}
非法输入应在进入求解器前返回结构化错误；不可行、无界、不收敛、时限、内存不足和后端缺失必须
相互区分。部分结果只有在对应 \fld{ValidityFlags}、\fld{model\_scope}、
\fld{model\_limitations} 或等价字段允许时才可使用。异常路径不得静默丢弃 AC/DC 资产、制造平衡
节点、把 NaN 改写为零或把近似解标为全局最优；系统原始对象应保持可恢复和可重试。

\subsection{与其他模块的接口关系}
接口链路遵循“工程语义模型 $\to$ 校验 $\to$ 投影/装配 $\to$ 求解或分析 $\to$ 结果回投影”。
跨模块传递时保留稳定 ID、单位、符号、时间基准和场景身份；消费潮流、OPF、动态或可靠性结果前，
先检查其收敛、覆盖范围和有效性标志。HTTP/JSON/Python 层只做语义保持适配，不重新解释求解结果。

\subsection{测试与验证建议}
最低验证矩阵包括：手算小例、标准算例、边界/异常输入、随机性质测试、算法或后端交叉校核、
序列化往返、稳定 ID 回投影、AC/DC 同号母线、重复运行确定性和规模扩展测试。数值算法还应加入
残差独立重算、雅可比有限差分或自动微分校核；近似模型应与高保真模型比较并给出误差分布，而非
只报告单个平均值。回归报告必须列明构建配置、算例版本、容差、通过/失败/跳过数及未验证范围。

\subsection{工业级评估指标}
\begin{itemize}
  \item \textbf{正确性}：守恒残差、约束最大违反、解析/基准误差、结果回投影一致性；
  \item \textbf{鲁棒性}：收敛率、不可行识别率、异常分类准确率、扰动敏感性与随机复现率；
  \item \textbf{性能}：端到端时延、吞吐量、峰值内存、并行效率与规模增长斜率；
  \item \textbf{可追溯性}：输入模式版本、稳定 ID 覆盖率、模型范围/限制字段完整率、告警可定位率；
  \item \textbf{工程有效性}：与标准、外部引擎或现场数据的偏差，以及误判/漏判对业务指标的影响。
\end{itemize}
指标应预先给出验收阈值和失败处置；未达到阈值时只能标记为实验性或受限适用。

\subsection{适用范围、局限性与后续改进}
本手册只评价当前源码覆盖的模型、后端和数据结构，不对未建模物理过程、未安装求解器或未执行的
外部验证作保证。后续改进应按“缺口 $\to$ 可量化指标 $\to$ 源码/测试变更 $\to$ 独立验证”闭环，
优先消除会造成错误结论的语义缺口，其次提升数值鲁棒性与性能，最后扩展模型范围。


\end{document}
